\itemhead{复习题二 · 第13题（教材p.62）}{XJ-B1-C02-RV-P062-Y1179}
\label{XJ-B1-C02-RV-P062-Y1179}
{\footnotesize 湘教版教材p.62 / PDF第69页 / 复习题二 / 第13题}
\field{原题}{如图，某新建居民小区欲建一面积为 $700\,\mathrm{m^2}$ 的矩形绿地，在绿地四周铺设人行道。设计要求绿地长边外人行道宽 $3\,\mathrm m$，短边外人行道宽 $4\,\mathrm m$。怎样设计绿地的长与宽，才能使人行道的占地面积最小（$\sqrt{21}\approx4.58$，结果精确到 $0.1\,\mathrm m$）？
\bankfig{C02-RV062}{0.32}}
\field{参考答案}{设长 $L$、宽 $W$，$LW=700$，人行道面积 $6L+8W+48$ 最小当 $6L=8W$。得 $L=20\sqrt{21}/3$ 米，$W=5\sqrt{21}$ 米。按题给 $\sqrt{21}\approx4.58$ 代入，约为30.5米、22.9米，与答案册一致；若用未截断的精确根值再舍入，长边为30.6米。保留精确式并注明采用的近似口径。}
\field{核对说明}{题给近似口径：按题给√21\ensuremath{\approx}4.58代入，长边约30.5米，与答案册PDF16一致；用精确根值再舍入则为30.6米。保留精确式并注明口径，不把近似输入差异当成答案册笔误。}
{\small \field{关联知识}{I03 基本不等式、等号与最值建模（教材帮p.52、53、54 / PDF61、62、63）；F03 函数求值、解析式与函数建模（教材帮p.80、87、88、89、90、131 / PDF89、96、97、98、99、140）}
\field{方法与易错点}{列正数条件与固定量，构造基本不等式，解等号条件并代回实际限制。 易错：把一个下界误当最小值；多个等号条件不能同时实现；漏单位和整数限制。}
\field{答案依据}{答案册相关题组：PDF第15、16页（册内页14、15）；按当前题库原题逐项重算，不按旧版题号机械对套。}}
